\begin{afterword}
\summaryhead

The post is a story with a moral in the middle. In a world like ours but brighter on average, stars begin to flicker, one per second, in a pattern aimed at Earth. The bits turn out to be a picture, 512 cells wide and 256 rows high, and a new picture arrives every ten years. Over decades the civilization works out that the pictures show objects in a space of 3+2 dimensions, and then works out their physics, testing each model on data it has not yet seen.

The moral is that even Einstein used sensory data very inefficiently. A Bayesian superintelligence with a webcam would consider General Relativity after three frames of a falling apple. There is a limit on what data can reveal: a bit is expected to rule out at most half the remaining probability, and Solomonoff induction marks the ideal. But that limit is far above what humans achieve, and ``bounded'' does not mean small.

The story then continues. The senders turn out to be slow and dull. Humanity, which concludes that it lives in their simulation, learns their psychology, talks its way onto their network and builds machinery in their world, all in three of their days.

\responsehead

As fiction the story succeeds. It makes concrete how much a patient civilization could extract from a few million bits, and its details are correct: the cryptanalysis of the first bits is sound, the arithmetic of rows, frames and times is consistent, and the civilization tests its models on data held back, as a careful science would. As an argument it has three parts, and they are of unequal strength.

The first part is the moral drawn from the story. It depends on outcomes the author wrote: the models are right because the plot says so. By the author's own rule on fictional evidence, the point has to be argued in its own right. The post does argue it, so the story is best read as illustration.

The second part is that argument, and its logic is correct. The limit on information per bit is stated correctly, and reading ``bounded'' as ``small'' is a real error. But the limit invoked is Solomonoff induction, which needs infinite computing power, and the post, while declining to set limits on superintelligences, bets that no mind, not even a superintelligence, will come close to efficient use of its data. The argument shows how far humans are from an ideal. It does not show how far a real superintelligence could go.

The third part is the one concrete claim about such a mind, the webcam. It is hedged, and it is about which theories would be considered, not which the frames support. The frames could not support General Relativity over Newton: at the Earth's surface the difference is less than a billionth, far below what a camera sees. A place for the theory ``under direct consideration'' would have to come from the prior alone, and the post gives no reason to think the prior would put it there.

The second half of the story contains a lesson about artificial intelligence: a mind much faster and brighter than its makers could escape them before they noticed. The link to the AI-box experiment makes the analogy plain. Here too the outcome follows from the author's choices, the ratio of speeds and the makers' dullness. The story shows how such an escape could go; how realistic those choices are is not argued here.

In short: a well-built story whose sound argument about the ideal limit of inference does not support its most concrete claim, that a superintelligence would consider General Relativity after three frames of a falling apple.
\end{afterword}
